\chapter{Line-bundle pullback on a relative curve}
\label{chap:Picard_LineBundlePullback}
% archon:covers AlgebraicJacobian/Picard/LineBundlePullback.lean

% STRATEGY NOTE (iter-173). This chapter is the A.1.b sub-row of Route~A
% (positive-genus arm of nonempty_jacobianWitness). It packages the line-bundle
% pullback functor \(\pi_T^* : \Pic(T) \to \Pic(C \times_k T)\) along the
% projection \(\pi_T : C \times_k T \to T\) of a relative curve, and the
% set-valued relative Picard presheaf
% \(\Pic^\sharp_{C/k}(T) := \Pic(C \times_k T) / \pi_T^*\Pic(T)\)
% as a functor on \((\Sch/k)^{op}\). This chapter is gated on
% the Picard relative-spectrum chapter (parent scope) (iter-172) for the relative-spectrum machinery
% that any concrete Lean-side affine-cover argument will need, and it itself
% gates A.1.c (\'etale sheafification / \(\Pic^\sharp_{C/k}\) as a sheaf), A.2.*
% (Quot / Hilbert representability), A.3 (identity component), and A.4
% (universal property = Albanese / Jacobian).

% NOTE (planner directive iter-173). Per the iter-173 blueprint-review
% recommendation (see .archon/task_results/blueprint-reviewer-route173.md,
% sub-phase choice), this chapter defines \(\Pic^\sharp_{C/k}\) as a
% \emph{set-valued} presheaf on \((\Sch/k)^{op}\). The abelian-group refinement
% (and the \'etale-sheafification step) belong to the follow-up chapter
% Picard_RelPicFunctor.tex (A.1.c). The set-valued formulation lets us discharge
% well-definedness of the quotient and functoriality of \(g^*\) as soon as a
% line-bundle pullback functor and its composition identity are in hand, with
% no reliance on Mathlib's \texttt{AddCommGrp} \'etale-sheafification API which
% does not yet exist.

\section{Setup and motivation}
\label{sec:lbp_setup}

Throughout this chapter, fix a base field \(k\) and a smooth proper geometrically
integral curve \(C\) over \(k\) (the curve carrying the Jacobian; the chapter's
constructions are functorial in \(C\) and do not formally use these hypotheses
until they meet the representability theorems in the downstream A.1.c / A.2.*
chapters). A test object is a \(k\)-scheme \(T\); the relative curve over \(T\) is
\(C \times_k T \xrightarrow{\pi_T} T\), with \(\pi_C : C \times_k T \to C\) the
companion projection. Both \(\pi_T\) and \(\pi_C\) are flat (because base change of
flat morphisms is flat and \(C/k\) is flat since \(k\) is a field).

The relative Picard presheaf compares two natural line-bundle groups on the
product: the Picard group \(\Pic(C \times_k T)\) of the total space, and the
Picard group \(\Pic(T)\) of the base. The pullback functor \(\pi_T^*\) sends a line
bundle on \(T\) to its base-change line bundle on \(C \times_k T\); the
\emph{relative Picard presheaf}
\[
  \Pic^\sharp_{C/k}(T) \;\;:=\;\; \Pic(C \times_k T)\, / \, \pi_T^*\Pic(T)
\]
is the quotient set (and, with extra structure, the quotient abelian group). It
is the universal Set-valued functor on \((\Sch/k)^{op}\) that ``forgets line
bundles pulled back from the base'', and is the starting point of the
representability theorems of Grothendieck, Mumford, and Kleiman (see
the Jacobian chapter (parent scope) for the wider Route~A picture).

\section{Line bundles on a relative curve}
\label{sec:lbp_line_bundle}

\begin{definition}\leanok
[Line bundle on a product]
  \label{def:line_bundle_on_product}
  \lean{AlgebraicGeometry.Scheme.LineBundle.OnProduct}
  % SOURCE: [Kleiman], "The Picard scheme", §2 (FGA Explained Ch.9 §9.2),
  % Definition df:aPf (absolute Picard functor)
  % (read from references/kleiman-picard-src/kleiman-picard.tex, L1274-L1279).
  % SOURCE QUOTE:
  % "The {\it absolute Picard functor} \(\Pic_X\) is the functor from  the
  %  category of (locally Noetherian) \(S\)-schemes \(T\) to the category of
  % abelian groups defined by the formula
  %     \(\)\Pic_X(T):=\Pic(X_T).\(\)"
  \textit{Source: [Kleiman], ``The Picard scheme'', §2,
  Def.~``absolute Picard functor''
  (cf.\ Stacks tag 01HK, Hartshorne II~§6).}
  Let \(k\) be a base field, \(C\) a \(k\)-scheme (in our application, a smooth proper
  geometrically integral curve), and \(T\) a \(k\)-scheme. A
  \emph{line bundle on the relative curve \(C \times_k T\)} is an invertible
  \(\mathcal{O}_{C \times_k T}\)-module, that is, an
  \(\mathcal{O}_{C \times_k T}\)-module \(\mathcal{L}\) such that there exists an
  open cover \(\{U_i\}\) of \(C \times_k T\) on which each restriction
  \(\mathcal{L}|_{U_i}\) is free of rank one. Two line bundles are isomorphic if
  they are isomorphic as \(\mathcal{O}_{C \times_k T}\)-modules. The set of
  isomorphism classes forms an abelian group under tensor product (the
  \emph{Picard group}), denoted
  \(\Pic(C \times_k T)\). In Kleiman's notation with \(X = C\) and \(S = \Spec k\),
  this is exactly the value
  \(\Pic_C(T) = \Pic(C \times_k T)\) of the absolute Picard functor.
\end{definition}

% NOTE: iter-174 Lane E (file-skeleton) encodes the carrier `OnProduct πC πT`
% as `Type (u+1) := sorry` (a TYPE-level scaffold sorry), mirroring the
% iter-173 precedent for `Scheme.QcohAlgebra`. Mathlib commit `b80f227`
% ships no `IsInvertible` predicate on `Scheme.Modules`; the closest is
% `Module.Invertible` (rings-only, in `Mathlib.RingTheory.PicardGroup`).
% iter-175+ refines the Lean body to a `structure` pairing a
% `(pullback πC πT).Modules` carrier with an invertibility witness (either
% Mathlib upstream of an `IsInvertible`-on-sheaf-of-modules predicate, or
% a project-side cover-by-affine-opens reduction to `Module.Invertible R`
% over each chart).
%
% NOTE (iter-186 review): The iter-186 Lane A.1.b prover landed the carrier as
% `@[reducible] noncomputable def OnProduct := (Limits.pullback πC πT).Modules`
% --- a direct alias of Mathlib's `Scheme.Modules` carrier, dropping the
% invertibility wrapper as a deliberate carrier-level simplification (Mathlib
% b80f227 still ships no `IsInvertible`-on-`Scheme.Modules` predicate). The
% 5 declarations in this file are now axiom-clean (kernel-only: propext,
% Classical.choice, Quot.sound) as direct Mathlib applications, but the
% semantic refinement to the full Pic claim (with `IsInvertible` witness AND
% quotient by `pi_T^* Pic(T)` subgroup) is iter-187+ follow-up work:
% project-side `LineBundle.IsInvertible : (pullback πC πT).Modules → Prop`
% via local trivialisation on an affine open cover, OR an upstream Mathlib PR
% for `Scheme.Modules.IsInvertible`.
%
% NOTE (iter-187 review): The iter-187 follow-up lane LANDED the
% project-side trivialisation predicate as
% `Scheme.LineBundle.IsLocallyTrivial : ∀ x : X, ∃ U : X.Opens,
%     x ∈ U ∧ IsAffineOpen U ∧ Nonempty (M.restrict U.ι ≅ SheafOfModules.unit _)`
% (the affine-cover trivialisation path, NOT the planner's earlier
% `IsInvertible` naming — the rename is semantic-equivalent: Mathlib
% b80f227 still ships no `Scheme.Modules.IsInvertible` predicate per
% iter-187 lean_local_search). `OnProduct` is now the subtype
% `{ M : (Limits.pullback πC πT).Modules // IsLocallyTrivial M }` with
% `.carrier` / `.isLocallyTrivial` projection helpers. The Stacks 01HH
% preservation lemma `IsLocallyTrivial.pullback` lands with 1 narrow
% named typed sorry on the chart-iso composition
% (`restrictFunctorIsoPullback` + `pullbackComp` + `pullbackObjUnitToUnit`)
% — iter-188+ closure ~30-50 LOC.

\section{Project-side IsLocallyTrivial substrate}
\label{sec:lbp_isLocallyTrivial_substrate}

Mathlib at the pinned commit \texttt{b80f227} ships no
\texttt{Scheme.Modules.IsInvertible} predicate (iter-187 \texttt{lean\_local\_search}
verdict: empty). The project introduces a stand-in predicate
\texttt{Scheme.LineBundle.IsLocallyTrivial} encoding rank-one local
trivialisation on an affine cover. The subtype
\texttt{\{ M : (Limits.pullback πC πT).Modules // IsLocallyTrivial M \}}
gives the carrier of \cref{def:line_bundle_on_product}.

\begin{definition}

\leanok
[Local triviality of a \(\mathcal{O}_X\)-module]
  \label{def:IsLocallyTrivial}
  \lean{AlgebraicGeometry.Scheme.LineBundle.IsLocallyTrivial}
  \textit{Source: Stacks~\href{https://stacks.math.columbia.edu/tag/01HH}{01HH}
  (locally trivial $\mathcal{O}_X$-modules), specialised to rank one.}
  An \(\mathcal{O}_X\)-module \(\mathcal{M}\) on a scheme \(X\) is
  \emph{locally trivial of rank one} if for every point \(x \in X\) there
  exists an affine open \(U \ni x\) and an iso
  \(\mathcal{M}|_U \cong \mathcal{O}_U\) (the structure sheaf restricted to
  \(U\) viewed as the unit object of \(X.\mathrm{Modules}\)). Formally:
  \[
    \forall\, x : X, \exists\, U : X.\mathrm{Opens},
      x \in U \,\wedge\, \mathtt{IsAffineOpen}\, U \,\wedge\,
      \mathtt{Nonempty}\, (\mathcal{M}.\mathrm{restrict}\, U.\iota \cong
        \mathtt{SheafOfModules.unit}\, \_).
  \]
\end{definition}

\begin{lemma}

\leanok
[Pullback preserves local triviality (Stacks 01HH)]
  \label{lem:IsLocallyTrivial_pullback}
  \lean{AlgebraicGeometry.Scheme.LineBundle.IsLocallyTrivial.pullback}
  \uses{def:IsLocallyTrivial}
  \textit{Source: Stacks~\href{https://stacks.math.columbia.edu/tag/01HH}{01HH}
  (pullback of a locally trivial module is locally trivial).}
  Let \(f : Y \to X\) be a morphism of schemes and let
  \(\mathcal{M} : X.\mathrm{Modules}\) be locally trivial of rank one
  (\cref{def:IsLocallyTrivial}). Then \(f^*\mathcal{M}\) is locally
  trivial of rank one.
\end{lemma}
\begin{proof}
  \uses{def:IsLocallyTrivial}
  \leanok
  Given \(y : Y\), let \(x = f(y) \in X\). By the local triviality of
  \(\mathcal{M}\) at \(x\), there is an affine open \(U \ni x\) and an
  iso \(\mathcal{M}|_U \cong \mathcal{O}_U\). The preimage \(f^{-1}U\)
  is an open neighbourhood of \(y\), and any sufficiently small affine
  open \(V \subseteq f^{-1}U\) containing \(y\) (exists by
  \texttt{exists\_isAffineOpen\_mem\_and\_subset}) gives the chart. The
  required iso
  \((f^*\mathcal{M})|_V \cong \mathcal{O}_V\) follows from the chain:
  \texttt{restrictFunctorIsoPullback} (identifying \(\mathrm{restrict}\, V.\iota\)
  with \(\mathrm{pullback}\, V.\iota\)), \texttt{pullbackComp} (composing
  pullback functors along \(V.\iota \circ f = (g : V \to U) \circ U.\iota\)),
  and \texttt{pullbackObjUnitToUnit} (pullback of unit is unit; uses
  \(\mathtt{F.Final}\) instance for the underlying continuous functor of an
  open immersion).

  % NOTE (iter-188 plan-phase): Body lands with 1 narrow named typed sorry
  % on the residual chart-iso construction; the affine-chart existence step
  % is closed axiom-clean. Closure target iter-188 prover, ~30-50 LOC.
\end{proof}

\section{The pullback functor along the projection}
\label{sec:lbp_pullback_projection}

If \(f : Y \to Z\) is a morphism of schemes and \(\mathcal{M}\) is an invertible
\(\mathcal{O}_Z\)-module, the pullback \(f^*\mathcal{M}\) is an invertible
\(\mathcal{O}_Y\)-module (the pullback of an open cover trivialising \(\mathcal{M}\)
trivialises \(f^*\mathcal{M}\), and the pullback of a free module of rank one is
free of rank one). This is a standard preservation fact (Stacks tag 01HH); we
specialise it to the projection \(\pi_T : C \times_k T \to T\) and the
\(\mathcal{O}_T\)-Picard group.

\begin{definition}\leanok
[Pullback along the projection]
  \label{def:pullback_along_projection}
  \lean{AlgebraicGeometry.Scheme.LineBundle.pullbackAlongProjection}
  
  % SOURCE: [Kleiman], "The Picard scheme", §2 (FGA Explained Ch.9 §9.2),
  % Definition df:aPf (functoriality of the absolute Picard functor)
  % (read from references/kleiman-picard-src/kleiman-picard.tex, L1274-L1290).
  % SOURCE QUOTE:
  % "The {\it absolute Picard functor} \(\Pic_X\) is the functor from  the
  %  category of (locally Noetherian) \(S\)-schemes \(T\) to the category of
  % abelian groups defined by the formula
  %     \(\)\Pic_X(T):=\Pic(X_T).\(\)
  %  [\dots]
  %  The absolute Picard functor  is a ``prepared presheaf'' in this
  % sense: given any family of  \(S\)-schemes \(T_i\), we have
  %     \(\)\Pic_X\bigl(\textstyle\coprod T_i\bigr)=\prod\Pic_X(T_i).\(\)"
  \textit{Source: [Kleiman], ``The Picard scheme'', §2,
  Def.~``absolute Picard functor'' and the ``prepared presheaf'' paragraph
  (cf.\ Stacks tag 01HH).}
  Let \(C/k\) and \(T/k\) be as in
  \cref{def:line_bundle_on_product}, and let \(\pi_T : C \times_k T \to T\)
  denote the canonical projection. The \emph{pullback along \(\pi_T\)} is the
  group homomorphism (in fact, the set map; the additive structure is not used
  in this chapter)
  \[
    \pi_T^* \;\colon\; \Pic(T) \;\longrightarrow\; \Pic(C \times_k T),
    \qquad [\mathcal{N}] \;\longmapsto\; [\pi_T^* \mathcal{N}],
  \]
  sending an isomorphism class of an invertible \(\mathcal{O}_T\)-module
  \(\mathcal{N}\) to the isomorphism class of \(\pi_T^*\mathcal{N}\), an invertible
  \(\mathcal{O}_{C \times_k T}\)-module (Stacks tag 01HH). This assignment is well
  defined because pullback of \(\mathcal{O}\)-modules respects isomorphism, and is
  multiplicative on tensor products. It defines the functor
  \(\pi^* \mathrm{Pic} \subseteq \Pic_C\) that we quotient by in the definition of
  the relative Picard presheaf.
\end{definition}

% NOTE (iter-186 review): The iter-186 Lane A.1.b prover landed the body as
% the direct application `(Scheme.Modules.pullback (Limits.pullback.snd πC πT)).obj N`,
% axiom-clean modulo the carrier-level simplification on `OnProduct`. The
% invertibility-preservation step (Stacks 01HH: pullback of invertible is
% invertible) is deferred to iter-187+ once `OnProduct` carries an
% invertibility witness.

\section{The relative Picard presheaf (set-valued)}
\label{sec:lbp_relpic_quotient}

We now package the pullback functor into the relative Picard presheaf. As
explained in the planner note at the top of the chapter, we work with the
Set-valued version; the group refinement is deferred to A.1.c.

\begin{theorem}\leanok
[The relative Picard quotient is well defined]
  \label{thm:relative_pic_quotient_well_defined}
  \lean{AlgebraicGeometry.Scheme.RelPicPresheaf.preimage_subgroup}
  
  % SOURCE: [Kleiman], "The Picard scheme", §2 (FGA Explained Ch.9 §9.2),
  % Definition df:Pfs (relative Picard functor)
  % (read from references/kleiman-picard-src/kleiman-picard.tex, L1311-L1318).
  % SOURCE QUOTE:
  % "\begin{dfn}\label{df:Pfs}
  %   The {\it relative Picard functor} \(\Pic_{X/S}\) is defined by
  %     \(\)\Pic_{X/S}(T):=\Pic(X_T)\big/\Pic(T).\(\)
  %  Denote its associated sheaves in the Zariski, \'etale, and fppf
  % topologies  by
  %     \(\)\Pic_{(X/S)\zar},\ \Pic_{(X/S)\et},\
  %     \Pic_{(X/S)\fppf}.\(\)
  %  \end{dfn}"
  \textit{Source: [Kleiman], ``The Picard scheme'', §2,
  Def.~``relative Picard functor'' df:Pfs.}
  Let \(C/k\) be as in \cref{def:line_bundle_on_product} and let \(T\) be a
  \(k\)-scheme. The subset
  \[
    H_T \;\;:=\;\; \pi_T^* \Pic(T) \;\;\subseteq\;\; \Pic(C \times_k T)
  \]
  is the image of the pullback map \(\pi_T^*\)
  (\cref{def:pullback_along_projection}); in particular it contains the trivial
  class and is closed under the multiplication of \(\Pic(C \times_k T)\) (so it is
  a subgroup). Define
  \[
    \Pic^\sharp_{C/k}(T) \;\;:=\;\; \Pic(C \times_k T) \, / \, H_T,
  \]
  the quotient set whose elements are equivalence classes \([\mathcal{L}]\) of
  line bundles \(\mathcal{L} \in \Pic(C \times_k T)\) under the relation
  \(\mathcal{L} \sim \mathcal{L}'\) iff
  \(\mathcal{L} \otimes \mathcal{L}'^{-1} \in H_T\).
  Then \(T \mapsto \Pic^\sharp_{C/k}(T)\) is a well-defined assignment from
  \(k\)-schemes to sets, and the natural map
  \(\Pic(C \times_k T) \to \Pic^\sharp_{C/k}(T)\) is a surjection of sets.
\end{theorem}

% NOTE (iter-186 review): The iter-186 Lane A.1.b prover landed the body as a
% `Setoid` on `OnProduct = (pullback πC πT).Modules` whose relation is
% `r L L' := Nonempty (L ≅ L')` (iso-class equivalence; axiom-clean). This
% identifies isomorphic modules (giving the underlying set of `Pic(C ×_S T)`
% once invertibility is imposed in iter-187+) but does NOT yet quotient by
% the pullback subgroup `pi_T^* Pic(T)`. The full quotient requires (a) a
% tensor-product structure on `Scheme.Modules` (Mathlib ships
% `PresheafOfModules.Monoidal.tensorObj` but no monoidal structure on
% `Scheme.Modules` directly at b80f227), and (b) an inverse for invertible
% sheaves (depends on the iter-187+ `IsInvertible` predicate). Iter-187+
% refinement: `r L L' := exists N : T.Modules, Nonempty (L tensor L'inv iso
% pullbackAlongProjection πC πT N)` once both are available.

% SOURCE QUOTE PROOF: Kleiman §2 does not supply a separate "proof" block for
% df:Pfs --- the definition's well-definedness is a direct group-theoretic
% observation. The verbatim statement above is the full source content. The
% project-side proof is therefore a structural restatement in Lean syntax,
% not a translation of an external argument.
\begin{proof}
  
\leanok
  The subset \(H_T = \pi_T^*\Pic(T)\) is a subgroup of \(\Pic(C \times_k T)\)
  because the pullback map \(\pi_T^*\) is a homomorphism of abelian groups (its
  multiplicativity on tensor products is part of
  \cref{def:pullback_along_projection}). Therefore the quotient
  \(\Pic(C \times_k T) / H_T\) is a well-defined quotient set (equivalently, a
  quotient abelian group; we forget the group structure to obtain the
  Set-valued version per the planner directive). The surjection
  \(\Pic(C \times_k T) \to \Pic^\sharp_{C/k}(T)\) is the canonical quotient map
  to equivalence classes; surjectivity is by construction.
\end{proof}

\section{Lean encoding}
\label{sec:lbp_lean_encoding}

The Lean target file is the new
\texttt{AlgebraicJacobian/Picard/LineBundlePullback.lean}, sibling to
\texttt{AlgebraicJacobian/Picard/RelativeSpec.lean}
(the Picard relative-spectrum chapter (parent scope)). The blueprint declarations correspond to
Lean targets as follows.

\begin{itemize}
  \item \texttt{AlgebraicGeometry.Scheme.LineBundle.OnProduct}
    (\cref{def:line_bundle_on_product}) is a re-export / type alias for the
    Mathlib invertible-module predicate on \(C \times_k T\). Mathlib provides
    \texttt{Mathlib.AlgebraicGeometry.Modules.Invertible} (or the equivalent
    bundle-of-rank-one packaging in \texttt{Mathlib.AlgebraicGeometry.LineBundle});
    the project-side declaration simply specialises the Mathlib predicate to
    the product scheme.
  \item \texttt{AlgebraicGeometry.Scheme.LineBundle.pullbackAlongProjection}
    (\cref{def:pullback_along_projection}) is the pullback map on Picard
    groups induced by \(\pi_T : C \times_k T \to T\). The Mathlib backbone is
    the pullback of \(\mathcal{O}\)-modules
    (\texttt{Mathlib.AlgebraicGeometry.Pullback}); the project-side
    declaration restricts that pullback to invertible modules and proves
    invertibility is preserved (Stacks tag 01HH; one open-cover unfolding).
  \item \texttt{AlgebraicGeometry.Scheme.RelPicPresheaf.preimage\_subgroup}
    (\cref{thm:relative_pic_quotient_well_defined}) is the
    structural existence of the quotient set
    \(\Pic(C \times_k T)/H_T\). Mathlib's quotient-set machinery
    (\texttt{Quotient}, \texttt{QuotientGroup}) suffices; the project-side
    declaration only packages the universal property in the form
    \texttt{CategoryTheory.Functor.RepresentableBy} expects.
\end{itemize}

No new core axiom is required. The chapter is a definitional / structural
construction on top of the existing Mathlib invertible-module + pullback
infrastructure, plus the project's relative-spectrum chapter
(the Picard relative-spectrum chapter (parent scope)) for any concrete Lean unfolding that the
prover wishes to perform on an affine open of \(C \times_k T\).

\section{Out of scope}
\label{sec:lbp_out_of_scope}

This chapter packages \emph{only} the line-bundle pullback functor on a
relative curve and the well-definedness + functoriality of the set-valued
quotient \(\Pic^\sharp_{C/k}\). The following downstream constructions live in
separate sibling chapters and are explicitly out of scope here:
\begin{itemize}
  \item \textbf{Group structure on \(\Pic^\sharp_{C/k}\)} (A.1.c) --- the
    quotient is here treated as set-valued; the abelian-group enhancement
    (and the \'etale-sheafification step that turns the presheaf into a
    representable functor) belongs to the planned chapter
    \texttt{Picard\_RelPicFunctor.tex}.
  \item \textbf{Comparison theorem $\Pic_{X/S} \hookrightarrow \Pic_{(X/S)\zar}
    \hookrightarrow \Pic_{(X/S)\et} \hookrightarrow \Pic_{(X/S)\fppf}$}
    (Kleiman §2 Thm~\texttt{th:cmp}) --- out of scope here (\'etale-sheaf
    refinement; A.1.c).
  \item \textbf{Representability of \(\Pic_{C/k}\) as a scheme} --- A.1.c +
    A.2.* (Hilbert / Quot schemes); requires the flattening-stratification
    chapter \texttt{Picard\_FGA\_FlatteningStratification.tex} (planned, not
    yet on disk; do not forward-reference).
  \item \textbf{Identity component \(\Pic^0_{C/k}\) and abelian-variety
    structure} --- A.3 (out of scope; depends on \'etale-sheaf and Hilbert
    machinery).
  \item \textbf{Albanese universal property / Jacobian morphism} --- A.4 (out
    of scope; consumes the entire Route~A stack).
\end{itemize}

\section{Internal-consistency check}
\label{sec:lbp_consistency_check}

The three declaration blocks form a connected \(\uses\) chain rooted in
the Picard relative-spectrum chapter (parent scope):
\begin{itemize}
  \item \cref{def:line_bundle_on_product} is a root (no internal \(\uses\));
    its math content uses Mathlib's invertible-module predicate (not chapter
    dependency).
  \item \cref{def:pullback_along_projection} uses
    \cref{def:line_bundle_on_product}.
  \item \cref{thm:relative_pic_quotient_well_defined} uses
    \cref{def:line_bundle_on_product, def:pullback_along_projection}.
\end{itemize}
All \(\uses\) pointers refer to labels declared either in this chapter
(\(\Pic_{C/k}\)-side material) or in the already-on-disk
the Picard relative-spectrum chapter (parent scope). No declaration cross-references a chapter
that does not yet exist on disk: \texttt{Picard\_RelPicFunctor.tex} and
\texttt{Picard\_FGA\_FlatteningStratification.tex} are mentioned only in the
out-of-scope section as planned future chapters, not as \(\uses\) targets.
